\documentclass[11pt,a4paper]{article}
\usepackage[T1]{fontenc}
\usepackage[utf8]{inputenc}
\usepackage{lmodern,microtype}
\usepackage[a4paper,margin=26mm,headheight=15pt]{geometry}
\usepackage{amsmath,amssymb,amsthm,mathtools}
\usepackage{booktabs,array,enumitem}
\usepackage[dvipsnames]{xcolor}
\usepackage{fancyhdr}
\usepackage[unicode,colorlinks=true,linkcolor=MidnightBlue,citecolor=MidnightBlue,urlcolor=MidnightBlue,breaklinks=true]{hyperref}
\usepackage{bookmark}
\usepackage{xurl}
\usepackage{listings}
\pagestyle{fancy}\fancyhf{}
\fancyhead[L]{\small Natural boundaries of quadratic feedback}
\fancyhead[R]{\small Research article}
\fancyfoot[C]{\thepage}
\renewcommand{\headrulewidth}{0.3pt}
\setlength{\parskip}{3pt}
\setlength{\emergencystretch}{3em}
\setlist{itemsep=4pt,topsep=5pt,leftmargin=1.8em}
\numberwithin{equation}{section}
\newtheorem{theorem}{Theorem}[section]
\newtheorem{proposition}[theorem]{Proposition}
\newtheorem{lemma}[theorem]{Lemma}
\newtheorem{corollary}[theorem]{Corollary}
\theoremstyle{definition}
\newtheorem{definition}[theorem]{Definition}
\newtheorem{example}[theorem]{Example}
\theoremstyle{remark}
\newtheorem{remark}[theorem]{Remark}
\newcommand{\C}{\mathbb C}
\newcommand{\R}{\mathbb R}
\newcommand{\N}{\mathbb N}
\newcommand{\ee}{\mathrm e}
\newcommand{\dd}{\mathrm d}
\newcommand{\Ree}{\operatorname{Re}}
\newcommand{\Imm}{\operatorname{Im}}
\newcommand{\dist}{\operatorname{dist}}
\newcommand{\B}{\mathcal B_0}
\newcommand{\U}{\widehat U}
\newcommand{\Q}{\widehat Q}
\newcommand{\inv}{\langle-1\rangle}
\newcommand{\TV}{\mathrm{TV}}
\newcommand{\coef}[2]{[#1^{#2}]}
\newcommand{\Dminus}{D_-}
\newcommand{\Hheat}{\mathcal H}
\newcommand{\res}{\operatorname{Res}}
\newcommand{\rationalamp}{\mathcal A}
\hypersetup{pdftitle={Natural Boundaries of Quadratic Exponential Feedback},pdfauthor={Research draft prepared with ChatGPT for Vladimir Reshetnikov},pdfsubject={Natural boundaries, quadratic exponential feedback, fine versus angular Borel summability, rational amplitudes}}
% ed. (2026-09-29): editorial-note environment for ProveIt cross-references;
% it uses no counter, so no author numbering changes.
\newenvironment{ednote}{\par\smallskip\begin{quote}\small\raggedright\textbf{Editorial note (ProveIt, 2026-09-29).}\ }{\end{quote}\smallskip}
\begin{document}
\begin{center}
{\LARGE\bfseries Natural Boundaries of\par Quadratic Exponential Feedback\par}
\vspace{0.7em}
{\large An exact obstruction to angular Borel summation,\par analytic-curve rigidity, and a rational-amplitude dichotomy\par}
\vspace{1em}
Research article prepared for Vladimir Reshetnikov\\
29 September 2026
\end{center}

\begin{abstract}
We study the formal feedback equation
\[
 \U(q)=\sum_{j\ge1}q^j\ee^{j^2\U(q)}
\]
and the literal left-side realization of its compositional inverse,
characterized by $\sum_{j\ge1}Q(u)^j\ee^{j^2u}=u$.
We prove that a neighborhood of zero in the imaginary axis is a natural
boundary for $Q$, although $Q$ is smooth up to that boundary.  Its integrated
Borel transform is entire and exponentially bounded in every fixed-width
neighborhood of the negative ray, but has no exponential bound in any open
angular neighborhood of that ray.  Consequently $\Q$, and also $\U$, is not
angularly $1$-summable in direction $\pi$.  This answers negatively the
specific open-arc question in the repository's negative-ray summation article.

The main tool is a local rigidity theorem for a heat expansion restricted
to an arbitrary analytic curve.  If the initial meromorphic function has
nearest poles at distances $s_\nu$, then the coefficient contribution of
a pole of order $p$ retains the nonzero multiplier
$\exp(-a'(0)s_\nu/2)$ under the curve $a=a(t)$.
Distinct squared pole distances prevent persistent cancellation.  At
rational imaginary times the quadratic kernel has rational generating data,
so an analytic implicit inverse would contradict this rigidity theorem.

We extend the argument to rational amplitude generating functions and
eventually quadratic real slopes.  Within this class, angular negative-direction
summability holds exactly for finite action support.  Nonpolynomial rational
amplitudes instead give a smooth natural boundary.  Finite-action inverses
nevertheless converge exponentially, with every fixed number of derivatives,
on a closed left half-disc.  Full proofs, exact-arithmetic checks,
120-digit numerical diagnostics, and further research questions are included.
\end{abstract}

\noindent\textbf{Status and scope.}
These are conventional mathematical proofs, not Lean-verified or externally
refereed results.  The proposed contribution is the implicit natural-boundary
theorem and its consequences for the stated feedback class.  Classical heat
summability, partial-theta methods, and factorial-composition mechanisms are
credited rather than claimed as new.  Global publication priority is not
established.  No general impossibility theorem for transseries summation is
asserted.

\clearpage
\setcounter{tocdepth}{1}
\tableofcontents
\clearpage

\section{The repository question and the distinction it requires}
\label{sec:context}

\subsection{The exact question answered}
The ProveIt repository contains a transseries development and a collection
of research articles on countable exponential feedback.  The snapshot
inspected for this article is
\begin{center}\small\ttfamily
0c973d8f5e7ef4550f4df1bf486d890037fd9f14.
\end{center}
The immediate source is \emph{Negative-Ray Summation of Countable
Exponential-Feedback Transseries} \cite{repo-negative}.  Its research agenda,
in the subsection on angular summability of the quadratic model, asks
whether an open arc of Borel directions containing $\pi$ can have the
uniform exponential-growth bounds required for angular summation.
The source explicitly distinguishes that problem from its proved
fixed-width-strip result.  The answer obtained here is \emph{no}.

The distinction is essential.  An entire Borel transform has no finite
Borel singularities, but can still grow too rapidly away from one ray to
permit rotation of the Laplace contour through any nonzero angle.
Conversely, an analytic function in a half-disc can be infinitely
differentiable on its diameter without admitting analytic continuation
through even one point of that diameter.  Our example realizes both
phenomena simultaneously, in a nonlinear implicit equation.

The related regularity and negative-ray articles already establish formal
existence, coefficient-growth information, and fine summation
\cite{repo-regularity,repo-negative}.  We do not claim those as new.
For the quadratic inverse we rederive the needed Borel representation and
its fine-summation property, with particularly simple frequency bounds.
This makes the main negative result independent of unverified asymptotic
claims in other research drafts.
% ed. (2026-09-29): relation to sibling packages this article could not see (filed after its snapshot).
\begin{ednote}
Packages filed in
\nolinkurl{Analysis/Transseries/docs/series-and-transseries/} after this
article's snapshot, and therefore not cited here, bear on it.
\nolinkurl{Sharp_Negative_Direction_Summability_Exponential_Feedback/}
proves sectorial $k$-summability in direction $\pi$ for every $k>1$ exactly
when $\lambda_j=O((j\log(j+1))^{1+1/k})$ and leaves the quadratic case
$k=1$ open; Theorem~\ref{thm:main} settles that case negatively in the
angular sense for $\lambda_j=j^2$. It thereby also answers, in the angular
sense, the regularity article's question on directional summability of the
quadratic model. The inverse coefficients of
\nolinkurl{Quadratic_Exponential_Feedback_After_Reversion/} and
\nolinkurl{Signed_Quadratic_Feedback_Inversion/} at $a=1$ agree with those
computed here through degree 16 (checked on filing, batch 49).
\end{ednote}

\subsection{What is not answered}
We do not determine the set of individual rays on which the Borel transform
has finite exponential type.  An isolated ray is not an open arc.  We do
not compute its precise growth in shrinking angular neighborhoods, the
best Taylor-truncation error, or nonlinear Stokes constants.  We also do
not infer differential transcendence from a natural boundary: nonlinear
differential equations can have solutions with natural boundaries.

Nor does the result contradict resurgence understood solely as suitable
analytic continuation of a Borel germ.  An entire Borel transform already
has unrestricted finite-plane continuation.  The obstruction here concerns
growth at infinity and the resulting sectorial Laplace sums, not the
existence of isolated finite Borel singularities.

\subsection{Relation to classical work}
The formal heat equation with meromorphic initial data is a classical
source of Gevrey expansions and Stokes phenomena.  Lutz, Miyake, and
Sch\"afke established foundational summability results for divergent heat
solutions \cite{lms}; Michalik and Podhajecka analyze Stokes phenomena for
meromorphic heat data \cite{mp}.  Han, Li, Sauzin, and Sun study partial
theta series, rational boundary points, discrete Fourier transforms, and
resurgence \cite{hlss}.  Sauzin distinguishes fine summation from summation
in an arc and proves nonlinear stability results \cite{sauzin}.

These works explain why heat poles and rational phases are natural tools.
The extra issue here is \emph{implicit cancellation}: the quantity playing
the role of the exponential amplitude is itself an unknown function of
time.  A fixed-amplitude gap theorem alone does not exclude the possibility
that this function removes a singularity.  Sections~\ref{sec:transfer}
and~\ref{sec:rigidity} prove the analytic-curve obstruction needed to
exclude precisely that possibility.

\section{Main statements and conventions}
\label{sec:main}

We use $\widehat f$ for a formal power series and $f$ for an analytic or
boundary-smooth realization.  The integrated Borel transform is
\begin{equation}
 \B\!\left(\sum_{n\ge0} f_n u^n\right)(\zeta)
       =\sum_{n\ge0}\frac{f_n}{n!}\zeta^n.
 \label{eq:borel-convention}
\end{equation}
Differentiating this transform gives the usual Borel minor of a series
with zero constant term.

\begin{definition}[Fine and angular negative-direction summability]
For $\delta>0$ write
\[
 T_\delta=\{\zeta:\dist(\zeta,(-\infty,0])<\delta\}.
\]
A series is \emph{finely summable in direction $\pi$} if its Borel germ
continues to some $T_\delta$ and there obeys
$|\B\widehat f(\zeta)|\le C\ee^{H|\zeta|}$.
It is \emph{angularly $1$-summable in direction $\pi$} if such a bound
holds in an open sector
$|\arg\zeta-\pi|<\eta$ for some $\eta>0$, together with a neighborhood
of the origin.  The argument is understood on a continuous determination
around $\pi$.  Equivalent definitions allow locally bounded exponential
constants as the direction varies; restricting to a smaller closed arc
recovers the uniform formulation used here.
\end{definition}

On the negative input ray the fine sum is
\begin{equation}
 \mathcal S_\pi\widehat f(-r)
   =\frac1r\int_0^\infty\ee^{-s/r}\B\widehat f(-s)\,\dd s,
 \qquad 0<r\ll1.
 \label{eq:negative-laplace}
\end{equation}
A fixed-width Borel neighborhood gives an input domain tangent to the
origin.  An angular Borel neighborhood permits rotated Laplace contours
and produces an input sector of opening strictly greater than $\pi$.
We give the precise contour argument used here in
Section~\ref{sec:angular}.

\begin{definition}[Natural boundary along an arc]
Let $f$ be holomorphic on one side of an open boundary arc $I$.  We say
that $I$ is a natural boundary for this branch of $f$ if no point of $I$
has a complex neighborhood on which a holomorphic function agrees with
$f$ on the original side.  This is a local statement about the specified
branch, not a claim about an unrelated function having the same formal
asymptotic expansion.
\end{definition}

Put
\begin{equation}
 K(q,u)=\sum_{j\ge1}q^j\ee^{j^2u},\qquad
 r_0=\frac1{32},\qquad
 D_- =\{u:|u|<r_0,\ \Ree u<0\}.
 \label{eq:kernel-domain}
\end{equation}
We use boundary values on $\Ree u=0$ whenever $|u|<r_0$.

\begin{theorem}[Quadratic feedback: smooth boundary, no angular sum]
\label{thm:main}
There is a unique small solution $Q$ of
\begin{equation}
 K(Q(u),u)=u,
 \qquad |Q(u)|\le\frac1{16},\quad
 |u|\le\frac1{32},\quad \Ree u\le0.
 \label{eq:inverse-equation}
\end{equation}
It is holomorphic in $D_-$, has continuous derivatives of every order on
the closed half-disc, and satisfies $Q(0)=0$, $Q'(0)=1$, and
\begin{equation}
 |Q'(u)-1|<\frac13.
 \label{eq:Qprime}
\end{equation}
The following conclusions hold.
\begin{enumerate}[label=\textup{(\roman*)}]
\item Every point of $\{it:|t|<r_0\}$ is a natural-boundary point for $Q$.
\item The formal Taylor series $\Q$ is the compositional inverse of the
unique formal solution $\U$ of
$\U(q)=K(q,\U(q))$.
Its integrated Borel transform is entire and exponentially bounded on
\emph{every} fixed-width $T_\delta$.  Its fine sum is $Q$.
\item For every $\eta>0$ and every $H>0$,
\begin{equation}
 \sup_{\substack{\zeta\ne0\cr |\arg\zeta-\pi|<\eta}}
       \ee^{-H|\zeta|}|\B\Q(\zeta)|=\infty.
 \label{eq:no-angular-bound}
\end{equation}
Neither $\Q$ nor $\U$ is angularly $1$-summable in direction $\pi$.
\item The inverse function $U=Q^{\inv}$ has the smooth arc
$\Gamma=\{Q(it):|t|<r_0\}$ as a natural boundary.  This arc is nowhere
real-analytic as a regular plane curve.
\end{enumerate}
\end{theorem}

The words ``every fixed-width'' do not imply a common exponential constant:
the constant can grow with $\delta$.  The impossibility in
\eqref{eq:no-angular-bound} concerns wedges of fixed positive angle,
whose width grows proportionally to their distance from zero.

The next theorem is a classification, not merely a perturbative example.

\begin{theorem}[Rational-amplitude dichotomy]
\label{thm:rational-main}
Let
\begin{equation}
 \rationalamp(z)=\sum_{j\ge1}c_jz^j
 \label{eq:amplitudes}
\end{equation}
be rational and holomorphic at zero, with $c_1\ne0$.  Suppose the real
slopes satisfy, for all sufficiently large $j$,
\begin{equation}
 \lambda_j=\alpha j^2+\beta j+\gamma,
 \qquad \alpha>0,
 \label{eq:eventual-quadratic}
\end{equation}
with arbitrary finite real modifications at smaller indices.  The small
left-side inverse defined by
\begin{equation}
 \sum_{j\ge1}c_jQ(u)^j\ee^{\lambda_j u}=u
 \label{eq:weighted-kernel}
\end{equation}
exists on a sufficiently small left half-disc and is smooth up to its
diameter.  Its formal series has an entire integrated Borel transform and
is finely summable in direction $\pi$, with sum $Q$.

The following conditions are equivalent:
\begin{enumerate}[label=\textup{(\alph*)}]
\item $\rationalamp$ is a polynomial, equivalently the action amplitudes
have finite support;
\item $Q$ extends holomorphically to a neighborhood of zero;
\item the formal series $\Q$ has positive radius of convergence;
\item $\Q$ is angularly $1$-summable in direction $\pi$.
\end{enumerate}
If these conditions fail, every point of a sufficiently small imaginary
interval is a natural-boundary point for $Q$.
\end{theorem}

The rationality hypothesis is substantive.  Rapidly decaying nonrational
amplitudes can make the kernel holomorphic across the boundary; an example
appears in Section~\ref{sec:examples}.  Thus the theorem does not replace a
joint amplitude--slope analysis for arbitrary amplitudes.

\section{The literal inverse on a closed half-disc}
\label{sec:inverse}

\subsection{Uniform contraction with explicit constants}
For $\Ree u\le0$ write \eqref{eq:inverse-equation} as
\begin{equation}
 q=T_u(q):=\ee^{-u}
       \left(u-\sum_{j\ge2}q^j\ee^{j^2u}\right).
 \label{eq:contraction}
\end{equation}
When $|q|\le1/16$ and $|u|\le1/32$,
\[
 |T_u(q)|\le\frac{32}{31}
       \left(\frac1{32}+\frac1{240}\right)
       =\frac{17}{465}<\frac1{16}.
\]
Here $\ee^{1/32}\le32/31$ follows from
$\ee^x\le(1-x)^{-1}$ for $0\le x<1$.
Moreover,
\begin{equation}
 |\partial_qT_u(q)|
  \le\frac{32}{31}
       \left(\frac1{(1-1/16)^2}-1\right)
  =\frac{32}{225}<1.
 \label{eq:contraction-constant}
\end{equation}
The contraction theorem gives the unique solution uniformly on the
closed parameter half-disc.  Iterating from zero gives holomorphic
functions on the open half-disc; uniform convergence proves holomorphy
of their limit.  Continuity up to the diameter follows from the same
uniform convergence.

\subsection{Smooth boundary jets and injectivity}
Every mixed derivative of $K$ is represented, on a slightly smaller
$q$-disc, by a uniformly convergent series of the form
\[
 \sum_{j\ge1}(j)_s j^{2t} q^{j-s}\ee^{j^2u},
 \qquad \Ree u\le0.
\]
Polynomial factors in $j$ are absorbed by the geometric decay in $q$.
The denominator $K_q$ stays bounded away from zero.  Implicit
differentiation, first in the interior and then by continuous extension,
therefore gives all boundary derivatives of $Q$.  Tangential derivatives
also follow directly by differentiating the uniformly contracting
boundary equation.  Repeating this argument proves that the one-sided
complex derivatives and the boundary tangential derivatives agree.
This is the precise meaning of smoothness on the closed half-disc.

For the quantitative first derivative, implicit differentiation gives
\begin{equation}
 Q'(u)=\frac{1-K_u(Q(u),u)}{K_q(Q(u),u)}.
 \label{eq:implicit-derivative}
\end{equation}
Set
\[
 d=\frac1{31}+\frac{31}{225},\qquad e=\frac{272}{3375}.
\]
The elementary geometric sums yield
\begin{align*}
 |K_q-1|&\le |\ee^u-1|+\sum_{j\ge2}j(1/16)^{j-1}\le d,\\
 |K_u|&\le\sum_{j\ge1}j^2(1/16)^j
     =\frac{(1/16)(1+1/16)}{(1-1/16)^3}=e.
\end{align*}
Since $(d+e)/(1-d)<1/3$, equation~\eqref{eq:implicit-derivative}
proves \eqref{eq:Qprime}.  Integrating $Q'-1$ on line segments in the
convex half-disc gives
\begin{equation}
 \frac23|u-v|\le |Q(u)-Q(v)|\le\frac43|u-v|.
 \label{eq:bilipschitz}
\end{equation}
The same inequality holds on the closure by continuity.  In particular,
$Q$ is univalent, its boundary parametrization is regular, and $Q(u)=0$
only when $u=0$.  The last fact also follows immediately from $K(0,u)=0$.

\subsection{Formal identification}
The operator
\[
 V\longmapsto\sum_{j\ge1}q^j\ee^{j^2V}
\]
improves agreement of series in $q\C[[q]]$ by one order.  It therefore has
a unique formal fixed point $\U$ with linear coefficient one.  Its formal
compositional inverse exists.  Substituting that inverse into the fixed
point equation gives the formal version of
\eqref{eq:inverse-equation}; conversely that equation determines its
coefficients successively.  The boundary derivatives just constructed
satisfy the same finite coefficient equations, so their Taylor series is
$\Q=\U^{\inv}$.

The beginning is
\begin{equation}
 \Q(u)=u-2u^2+\frac12u^3-\frac23u^4
       -\frac{29}{8}u^5-\frac{743}{60}u^6
       -\frac{8491}{144}u^7-\frac{92726}{315}u^8+\cdots.
 \label{eq:first-coefficients}
\end{equation}
This identity is formal; it does not assert convergence.

\section{Factorial expansions along analytic curves}
\label{sec:transfer}

We now prove the estimate that prevents an unknown analytic amplitude
from hiding a heat singularity.  The proof is elementary and includes
the uniform bound needed to justify the limiting coefficient sum.

For an integer $p\ge1$ define
\begin{equation}
 D_{p,n}=\frac{\Gamma(2n+p)}{\Gamma(p)n!},\qquad
 F_p(z)=\sum_{n\ge0}D_{p,n}z^n.
 \label{eq:Dpn}
\end{equation}
The second expression is a formal series.  Differentiating a pole of
order $p$ twice at each heat step produces exactly these coefficients.

\begin{lemma}[Analytic composition of a factorial series]
\label{lem:composition}
Let $p\ge1$ be an integer, let $s(t)$ be analytic near zero with
$s(0)=s_0\ne0$, and put $s_1=s'(0)$.  Then
\begin{equation}
 \coef{t}{n}\left\{
  \sum_{\ell\ge0}\frac{D_{p,\ell}t^\ell}{s(t)^{2\ell+p}}
 \right\}
 =D_{p,n}s_0^{-2n-p}
       \left(\ee^{-s_0s_1/2}+o(1)\right).
 \label{eq:pole-transfer}
\end{equation}
The $o(1)$ is additive inside the displayed parentheses.  In particular,
it does not assert a relative error for a sum of several pole
contributions that might be small by oscillation.
\end{lemma}

\begin{proof}
Let $R(t)=s_0^2/s(t)^2$ and $h(t)=s(t)^{-p}$.  Thus $R(0)=1$ and
$R'(0)=-2s_1/s_0$.  If $b_n$ is the coefficient on the left of
\eqref{eq:pole-transfer}, then the finite coefficient identity
\begin{equation}
 \frac{b_n}{D_{p,n}s_0^{-2n}}
  =\sum_{k=0}^n
       \frac{D_{p,n-k}}{D_{p,n}}s_0^{2k}
       \coef{t}{k}\{h(t)R(t)^{n-k}\}
 \label{eq:composition-sum}
\end{equation}
follows by setting $k=n-\ell$.

For fixed $k$, the exact ratio
\[
 \frac{D_{p,m+1}}{D_{p,m}}
       =\frac{(2m+p)(2m+p+1)}{m+1}
\]
gives
$D_{p,n-k}/D_{p,n}\sim(4n)^{-k}$.
Since $R(0)=1$, the coefficient
$\coef{t}{k}\{h(t)R(t)^{n-k}\}$ is a polynomial in $n$ of degree at
most $k$, with leading coefficient $h(0)R'(0)^k/k!$.  The $k$th summand
in \eqref{eq:composition-sum} therefore tends to
\[
 h(0)\frac{(s_0^2R'(0)/4)^k}{k!}.
\]

To justify summation of these limits, first note that, for all $m\ge0$
and $p\ge1$,
\[
 \frac{D_{p,m+1}}{D_{p,m}}\ge2(m+1).
\]
The geometric mean of the largest $k$ integers in $\{1,\ldots,n\}$
is at least the geometric mean of all $n$ of them, and
$n!\ge(n/\ee)^n$.  Consequently, for $1\le k\le n$,
\begin{equation}
 \frac{D_{p,n-k}}{D_{p,n}}
 \le2^{-k}\frac{(n-k)!}{n!}
 \le\left(\frac{\ee}{2n}\right)^k.
 \label{eq:ratio-majorant}
\end{equation}
Choose $\delta>0$ small enough that $R$ has an analytic logarithm on
$|t|\le\delta$.  There are constants $H,M$ such that
$|h(t)|\le H$ and $|R(t)|\le\ee^{M|t|}$ there.
Choose $B\ge\max(M,1/\delta,1)$.
Cauchy's coefficient estimate on the circle $|t|=k/(Bn)$ gives
\[
 \left|\coef{t}{k}\{h(t)R(t)^{n-k}\}\right|
       \le H\left(\frac{\ee Bn}{k}\right)^k.
\]
Combined with \eqref{eq:ratio-majorant}, the absolute value of the
$k$th summand of \eqref{eq:composition-sum} is at most
\begin{equation}
 H\left(\frac{|s_0|^2\ee^2B}{2k}\right)^k.
 \label{eq:summable-majorant}
\end{equation}
This majorant is summable over $k\ge1$ and independent of $n$.
The term $k=0$ is bounded separately.  Dominated convergence in the
counting measure now proves that \eqref{eq:composition-sum} tends to
\[
 h(0)\exp\left(\frac{s_0^2R'(0)}4\right)
       =s_0^{-p}\ee^{-s_0s_1/2}.
\]
This is \eqref{eq:pole-transfer}.
\end{proof}

\begin{remark}[Why the first derivative, but not higher derivatives, survives]
For a fixed loss of $k$ heat indices, the factorial ratio contributes
$n^{-k}$ and the coefficient of the analytic substitution contributes
$n^k$.  Only repeated uses of the linear coefficient of $R(t)-1$ reach
that highest degree in $n$.  Higher derivatives of $s(t)$ affect lower
orders.  The exponential multiplier is never zero, which is the feature
needed for a noncancellation theorem.
\end{remark}

\begin{lemma}[Analytic remainder estimate]
\label{lem:analytic-remainder}
Let $\psi$ be analytic on $|a-a_0|<R_*$ and let $a(t)$ be analytic near
zero with $a(0)=a_0$.  For every $R_1<R_*$ there is a constant $C$ such
that, for all sufficiently large $n$,
\begin{equation}
 \left|\coef{t}{n}
   \sum_{\ell\ge0}\frac{t^\ell}{\ell!}
               \psi^{(2\ell)}(a(t))\right|
       \le C D_{1,n}R_1^{-2n}.
 \label{eq:analytic-remainder}
\end{equation}
\end{lemma}

\begin{proof}
Choose $R_2$ with $R_1<R_2<R_*$.  If
$\psi(a_0+x)=\sum c_mx^m$, Cauchy's estimate gives
$|c_m|\le M R_2^{-m}$.  Write
$a(t)-a_0=\sum_{j\ge1}a_jt^j$ and
$A(t)=\sum_{j\ge1}|a_j|t^j$.
Coefficientwise absolute values of the series in
\eqref{eq:analytic-remainder} are bounded by the coefficients of
\[
 M R_2\sum_{\ell\ge0}
       \frac{D_{1,\ell}t^\ell}{(R_2-A(t))^{2\ell+1}}.
\]
Indeed $M R_2/(R_2-x)$ majorizes $\psi(a_0+x)$, and taking an even
number of derivatives and substituting $A$ preserve this majorization.
Lemma~\ref{lem:composition} bounds these coefficients by
$C'D_{1,n}R_2^{-2n}$ for large $n$.  This is stronger than
\eqref{eq:analytic-remainder}.
\end{proof}

\section{A heat-pole rigidity theorem}
\label{sec:rigidity}

Let $\phi$ be meromorphic near a closed disc centered at $a_0$, analytic
at $a_0$, and with at least one pole.  Let $R>0$ be the distance from
$a_0$ to the nearest poles, and denote those poles by
$\omega_1,\ldots,\omega_d$.  Put
\begin{equation}
 s_\nu=a_0-\omega_\nu,\qquad |s_\nu|=R.
 \label{eq:pole-distances}
\end{equation}
Suppose $\phi$ has no other singularity in some larger disc
$|a-a_0|<R_*$, $R_*>R$.  Let $p$ be the largest pole order among the
nearest poles, and let $J$ index those of order $p$.  For $\nu\in J$
write the highest Laurent term as
$r_\nu(a-\omega_\nu)^{-p}$, with $r_\nu\ne0$.

\begin{theorem}[Rigidity under analytic curves]
\label{thm:rigidity}
For every analytic $a(t)$ with $a(0)=a_0$, the formal series
\begin{equation}
 \widehat G(t)=\sum_{\ell\ge0}\frac{t^\ell}{\ell!}
                    \phi^{(2\ell)}(a(t))
       =\sum_{n\ge0}g_nt^n
 \label{eq:curved-heat}
\end{equation}
satisfies
\begin{equation}
 g_n=D_{p,n}R^{-2n}
   \left\{\sum_{\nu\in J}
       r_\nu s_\nu^{-p}\ee^{-a'(0)s_\nu/2}
       \left(\frac{R}{s_\nu}\right)^{2n}+o(1)\right\}.
 \label{eq:curved-pole-asymptotic}
\end{equation}
If the numbers $s_\nu^2$, $\nu\in J$, are pairwise distinct, then
\begin{equation}
 \limsup_{n\to\infty}
      \left(\frac{|g_n|}{n!}\right)^{1/n}
       =\frac4{R^2}>0.
 \label{eq:exact-local-type}
\end{equation}
In particular $\widehat G$ cannot be the Taylor series of a holomorphic
germ.  Pairwise distinctness holds automatically if all the nearest
poles lie strictly to the left of $a_0$, that is,
$\Ree(a_0-\omega_\nu)>0$.
\end{theorem}

\begin{proof}
Subtract all principal parts at the nearest poles.  The remainder is
analytic on a disc of radius greater than $R$.
For a principal part $r(a-\omega)^{-h}$, its contribution to
\eqref{eq:curved-heat} is
\[
 r\sum_{\ell\ge0}
       \frac{D_{h,\ell}t^\ell}{(a(t)-\omega)^{2\ell+h}}.
\]
Lemma~\ref{lem:composition} gives its coefficient asymptotic.  If
$h<p$, then
\[
 \frac{D_{h,n}}{D_{p,n}}
       \sim\frac{\Gamma(p)}{\Gamma(h)}(2n)^{h-p}\longrightarrow0.
\]
Lemma~\ref{lem:analytic-remainder}, with $R<R_1<R_*$, makes the analytic
remainder exponentially smaller than $D_{p,n}R^{-2n}$.
This proves \eqref{eq:curved-pole-asymptotic}.

Set
\[
 A_\nu=r_\nu s_\nu^{-p}\ee^{-a'(0)s_\nu/2},\qquad
 \rho_\nu=(R/s_\nu)^2.
\]
All $A_\nu$ are nonzero and all $|\rho_\nu|=1$.  Distinct squared
distances mean distinct $\rho_\nu$.  The finite geometric-sum identity
implies
\begin{equation}
 \frac1N\sum_{n=0}^{N-1}
       \left|\sum_{\nu\in J}A_\nu\rho_\nu^n\right|^2
     \longrightarrow\sum_{\nu\in J}|A_\nu|^2>0.
 \label{eq:cesaro}
\end{equation}
Thus the expression in braces in
\eqref{eq:curved-pole-asymptotic} is bounded and bounded away from zero
along an infinite subsequence.  Finally,
\[
 \left(\frac{D_{p,n}}{n!}\right)^{1/n}\longrightarrow4
\]
by Stirling's formula, or by the central-binomial coefficient estimate
and the finite product relating $D_{p,n}$ to $D_{1,n}$.
Equation~\eqref{eq:exact-local-type} follows.

If $s_\nu^2=s_\mu^2$ with $\nu\ne\mu$, then
$s_\nu=-s_\mu$.  This is impossible when both have strictly positive
real part.  A holomorphic germ has exponentially bounded Taylor
coefficients and hence zero value of the limsup in
\eqref{eq:exact-local-type}, proving the final assertion.
\end{proof}

\subsection{From a formal heat expression to actual boundary derivatives}
Suppose $f_j$ is periodic and $a(t)$ is analytic with $\Ree a(0)>0$.
For $\Ree t\le0$ and sufficiently small $t$, define
\begin{equation}
 G(t)=\sum_{j\ge1}f_j\ee^{-ja(t)+j^2t},\qquad
 \phi(a)=\sum_{j\ge1}f_j\ee^{-ja}.
 \label{eq:actual-curved-heat}
\end{equation}
Every fixed derivative is uniformly dominated by a polynomial in $j$
times $\ee^{-cj}$ for some $c>0$.  Differentiation to the boundary is
therefore legitimate.  Expanding a fixed derivative, or collecting the
finite Taylor coefficient of each exponential, shows that the Taylor
jet of $G$ is exactly \eqref{eq:curved-heat}.  In particular, if $G$ had
a holomorphic continuation through zero, that continuation would have
those coefficients.  The same argument applies to exponentially
bounded $f_j$ whenever $\Ree a(0)$ is sufficiently large.

This observation does not replace the divergent series by a convergent
one.  It only identifies each finite derivative, which is all the
contradiction requires.

\section{Rational phases force a natural boundary}
\label{sec:natural}

\subsection{The finite Fourier data}
Fix a nonzero rational imaginary point
\begin{equation}
 u_0=\frac{2\pi i r}{m},\qquad r\in\mathbb Z,\quad m\in\N_{>0},
 \qquad |u_0|<r_0.
 \label{eq:rational-time}
\end{equation}
The sequence $f_j=\ee^{u_0j^2}$ is $m$-periodic, since
$(j+m)^2-j^2=2jm+m^2$.
Its generating function is rational in $z=\ee^{-a}$:
\begin{equation}
 \phi(a)=\sum_{j\ge1}f_j\ee^{-ja}
     =\frac{\sum_{j=1}^m f_j\ee^{-ja}}{1-\ee^{-ma}}.
 \label{eq:periodic-rational}
\end{equation}
Equivalently, with the finite Fourier coefficients
\begin{equation}
 \rho_k=\frac1m\sum_{j=1}^m
             \exp\left(\frac{2\pi i(rj^2-kj)}m\right),
       \qquad 0\le k<m,
 \label{eq:gauss-fourier}
\end{equation}
we have
\begin{equation}
 \phi(a)=\sum_{k=0}^{m-1}
       \frac{\rho_k}{\ee^{a-2\pi ik/m}-1}.
 \label{eq:fourier-poles}
\end{equation}
Its poles are simple, lie on the imaginary axis, and have residues
$\rho_k$ at $2\pi ik/m+2\pi i\mathbb Z$ whenever $\rho_k\ne0$.
At least one such residue is nonzero, because the finite Fourier
transform is invertible and $f_j$ is not the zero sequence.
No nonvanishing theorem for a particular quadratic Gauss sum is needed.

\subsection{Why an analytic implicit inverse is impossible}
Suppose, for contradiction, that $Q$ extends holomorphically across $u_0$.
Since $Q(u_0)\ne0$ and $|Q(u_0)|<1$, a local analytic logarithm gives
\begin{equation}
 a(t)=-\log Q(u_0+t),\qquad \Ree a(0)>0.
 \label{eq:analytic-amplitude}
\end{equation}
On the left side of $t=0$, the defining equation becomes
\begin{equation}
 \sum_{j\ge1}f_j\ee^{-ja(t)+j^2t}=u_0+t.
 \label{eq:contradiction-identity}
\end{equation}
The Taylor coefficients of the left side are those of the curved heat
series generated by \eqref{eq:periodic-rational}.
All its poles lie strictly to the left of $a(0)$ in real part.
Theorem~\ref{thm:rigidity} therefore gives a strictly positive
factorial-normalized type.  The right side has no coefficient of degree
at least two.  This contradiction proves that $Q$ has no holomorphic
continuation at $u_0$.

Points \eqref{eq:rational-time} with $u_0\ne0$ are dense in the imaginary
interval $|\Imm u|<r_0$.  A continuation across any point of that interval
would give a continuation across nearby rational points.  Hence every
point, including zero, is a natural-boundary point.  This proves
Theorem~\ref{thm:main}(i).

\begin{remark}[Why a fixed-amplitude argument is insufficient]
For a fixed $q$, the series $\sum q^j\ee^{j^2u}$ has a familiar lacunary
structure after putting $z=\ee^u$.  But the implicit identity uses
$q=Q(u)$, not a constant $q$.  The nonzero multiplier in
\eqref{eq:curved-pole-asymptotic} is what excludes cancellation by an
analytic $Q$.  This is the additional argument needed to pass from a
partial theta function to an implicitly defined branch.
\end{remark}

\subsection{The natural boundary of the forward realization}
By \eqref{eq:bilipschitz}, $Q$ maps $D_-$ conformally onto
$\Omega=Q(D_-)$ and maps its diameter homeomorphically onto $\Gamma$.
Let $U=Q^{\inv}$ on $\Omega$.  Suppose $U$ had a holomorphic
continuation across $q_0=Q(it_0)$.
Its derivative there would be the interior limit
$1/Q'(it_0)\ne0$.  The holomorphic inverse function theorem would then
produce a local holomorphic inverse of this continuation near $it_0$.
On the original side that inverse equals $Q$, contradicting the natural
boundary just proved.

If $\Gamma$ were a regular real-analytic arc near any of its points,
an analytic parametrization with nonzero derivative would extend to a
local holomorphic coordinate straightening the arc to a real interval.
In that coordinate $Q$ has real boundary values on a straight interval.
Schwarz reflection extends it holomorphically across that interval;
undoing the coordinate contradicts the same natural boundary.
Thus $\Gamma$ is nowhere real-analytic as a regular curve.

Its smooth jet at the origin nevertheless begins
\begin{equation}
 Q(it)=it+2t^2-\frac{i}{2}t^3-\frac23t^4+O(t^5).
 \label{eq:boundary-jet}
\end{equation}
A low-order polynomial plot can therefore look completely regular and
almost parabolic.  The obstruction is not visible in a finite jet.

\section{An entire Borel transform and its fine sum}
\label{sec:borel}

This section reproduces, in the quadratic specialization, the required
summation facts.  The signed Lagrange-block method comes from the
repository's negative-ray article \cite{repo-negative}; the simplified
frequency bounds below make the present proof short.

\subsection{Exact exponential blocks}
Invert $q\mapsto K(q,u)$ while keeping $u$ temporarily fixed and the
target $w$ independent.  The coefficient form of Lagrange inversion
\cite{gessel} gives
\[
 q(w;u)=\sum_{k\ge1}\frac{w^k}{k}
        \coef{q}{k-1}\left(\sum_{j\ge1}q^{j-1}\ee^{j^2u}\right)^{-k}.
\]
Setting $w=u$ is formally legitimate.  For $k\ge2$ and an ordered
composition $\mathbf m=(m_1,\ldots,m_\ell)$ of $k-1$ into positive
parts, define
\begin{equation}
 w_{k,\ell}=\frac{(-1)^\ell}{k}\binom{k+\ell-1}{\ell},\qquad
 b_{k,\mathbf m}=\sum_{i=1}^{\ell}m_i^2+k-2.
 \label{eq:block-weights}
\end{equation}
The exact block representation is
\begin{equation}
 \Q(u)=u\ee^{-u}
       +\sum_{k\ge2}u^k
          \sum_{\mathbf m\vDash k-1}
                    w_{k,\ell(\mathbf m)}\ee^{b_{k,\mathbf m}u}.
 \label{eq:quadratic-blocks}
\end{equation}
Indeed the exponent before simplification is
$\sum_i(m_i+1)^2-(k+\ell)$, which equals the expression in
\eqref{eq:block-weights}.  The important bounds are
\begin{equation}
 2k-3\le b_{k,\mathbf m}\le k^2-k-1<k^2,
 \qquad k\ge2.
 \label{eq:frequency-range}
\end{equation}
Thus the only negative-frequency block is the initial $u\ee^{-u}$.

The uncancelled total mass of the $k$th block is
\begin{equation}
 C_k=\frac1k\sum_{\ell=1}^{k-1}
       \binom{k+\ell-1}{\ell}\binom{k-2}{\ell-1},\qquad C_1=1.
 \label{eq:mass}
\end{equation}
A second application of Lagrange inversion, or substitution in the
quadratic equation, yields
\begin{equation}
 \sum_{k\ge1}C_kt^k
   =\frac{1+t-\sqrt{1-6t+t^2}}4.
 \label{eq:mass-gf}
\end{equation}
The coefficients are nonnegative and the value at $t=1/6$ is $1/4$.
Consequently
\begin{equation}
 C_k\le\frac{6^k}{4},\qquad k\ge1.
 \label{eq:mass-bound}
\end{equation}
For $\Ree u\le0$, the blocks with $k\ge2$ have modulus at most
$C_k|u|^k$.  Hence \eqref{eq:quadratic-blocks} converges normally on a
small left half-disc and equals the literal inverse there, by ordinary
Lagrange inversion in $q$ and the uniqueness from
Section~\ref{sec:inverse}.

\subsection{A normalized Bessel kernel}
For $k\ge1$ and real $b$ put
\begin{equation}
 B_{k,b}(\zeta)=\sum_{m\ge0}
      \frac{b^m\zeta^{m+k}}{m!(m+k)!}.
 \label{eq:one-borel-block}
\end{equation}
This is the integrated Borel transform of $u^k\ee^{bu}$.
The probability density
\[
 p_k(v)=\frac{\Gamma(k+1)}{\sqrt\pi\,\Gamma(k+1/2)}
             (1-v^2)^{k-1/2},\qquad -1<v<1,
\]
has even moments
\[
 \int_{-1}^1v^{2m}p_k(v)\,\dd v
       =\frac{(2m)!k!}{4^m m!(m+k)!}.
\]
Expanding the exponential and using those moments proves
\begin{equation}
 B_{k,b}(\zeta)=\frac{\zeta^k}{k!}
        \int_{-1}^1\ee^{2\sqrt{b\zeta}\,v}p_k(v)\,\dd v.
 \label{eq:bessel-integral}
\end{equation}
The integral is independent of the square-root choice.  In particular,
for $b\ge0$,
\begin{equation}
 |B_{k,b}(\zeta)|\le\frac{|\zeta|^k}{k!}
          \ee^{2\sqrt b\,|\Ree\sqrt\zeta|},\qquad
 |B_{k,b}(-s)|\le\frac{s^k}{k!}.
 \label{eq:bessel-bounds}
\end{equation}
For $b<0$, the same integral gives the less sharp but sufficient bound
$|B_{k,b}(\zeta)|\le |\zeta|^k\ee^{2\sqrt{|b||\zeta|}}/k!$.

\subsection{Entire continuation and tube bounds}
Apply \eqref{eq:one-borel-block} termwise to
\eqref{eq:quadratic-blocks}.  On $|\zeta|\le R$, equations
\eqref{eq:frequency-range} and \eqref{eq:mass-bound} bound the $k$th
block by
\[
 \frac{6^k}{4}\frac{R^k}{k!}\ee^{2k\sqrt R}.
\]
The sum converges for every $R$, so the resulting Borel function is entire.
Its Taylor coefficients at zero equal those of $\B\Q$, because each
ordinary coefficient uses only finitely many outer indices $k$.

For $\zeta\in T_\delta$,
\begin{equation}
 |\Ree\sqrt\zeta|^2
       =\frac{|\zeta|+\Ree\zeta}{2}\le\delta.
 \label{eq:sqrt-tube}
\end{equation}
This follows from the $2$-Lipschitz property of
$\zeta\mapsto|\zeta|+\Ree\zeta$, which vanishes on the negative ray.
Therefore
\begin{equation}
 |\B\Q(\zeta)|
 \le |\zeta|\ee^{2\sqrt{|\zeta|}}
      +\frac14\exp\left(6\ee^{2\sqrt\delta}|\zeta|\right),
 \qquad \zeta\in T_\delta.
 \label{eq:explicit-tube-bound}
\end{equation}
The first term is also bounded by a constant times an exponential in
$|\zeta|$.  This proves fine summability on every fixed-width tube.

\subsection{Identification of the fine sum}
For each fixed $k,b$, integrating the entire series in
\eqref{eq:one-borel-block} termwise against $r^{-1}\ee^{-s/r}$ is
justified by the integrable absolute series.  It gives
\begin{equation}
 \frac1r\int_0^\infty\ee^{-s/r}B_{k,b}(-s)\,\dd s
       =(-r)^k\ee^{-br}.
 \label{eq:block-laplace}
\end{equation}
To exchange the \emph{infinite block sum} with the integral, use the
stronger ray estimate in \eqref{eq:bessel-bounds}:
\[
 \sum_{k\ge2}\sum_{\mathbf m\vDash k-1}
     |w_{k,\ell}B_{k,b_{k,\mathbf m}}(-s)|
       \le\sum_{k\ge2}\frac{6^ks^k}{4k!}\le\frac14\ee^{6s}.
\]
The initial block is bounded by $s\ee^{2\sqrt s}$.
Thus dominated convergence applies when $0<r<1/6$, after also restricting
$r$ to the small inverse domain.  Equation~\eqref{eq:block-laplace}
then returns exactly \eqref{eq:quadratic-blocks} at $u=-r$.
This proves $\mathcal S_\pi\Q(-r)=Q(-r)$.
The same identity holds on the common complex domain by holomorphy.

\section{Why no angular neighborhood can work}
\label{sec:angular}

\begin{lemma}[Angular growth produces a larger input sector]
\label{lem:sector-laplace}
Let $B$ be holomorphic near zero and on
$|\arg\zeta-\pi|<\eta$, and suppose
$|B(\zeta)|\le C\ee^{H|\zeta|}$ on a smaller closed angular sector
still containing $\pi$.  Its Laplace transforms along those rays glue
to a holomorphic function in an input sector centered at $\pi$ with
opening strictly greater than $\pi$.  On the negative input ray it is
the integral \eqref{eq:negative-laplace}.
\end{lemma}

\begin{proof}
Take $0<\eta_1<\eta$ on which the exponential bound is uniform.  For a
ray of direction $\theta\in[\pi-\eta_1,\pi+\eta_1]$, the integral
\[
 \frac1u\int_0^{\ee^{i\theta}\infty}
             \ee^{-\zeta/u}B(\zeta)\,\dd\zeta
\]
converges whenever $\Ree(\ee^{i\theta}/u)>H$.
On overlaps, contour deformation between two sufficiently close rays
is justified by the exponential bound: the large joining arc has
exponentially decreasing integrand, and the small arc tends to zero.
The integrals therefore agree and glue.
Choose $0<\eta_2<\eta_1$.  Every direction
$|\arg u-\pi|<\pi/2+\eta_2$ makes an angle strictly less than $\pi/2$
with an allowed Laplace ray; after decreasing the input radius uniformly
on a slightly smaller closed sector, the convergence condition holds.
The resulting opening is $\pi+2\eta_2>\pi$.
\end{proof}

\begin{proof}[Proof of Theorem~\ref{thm:main}(iii)]
Suppose an exponential bound existed for $\B\Q$ in an angular
neighborhood of $\pi$.  Lemma~\ref{lem:sector-laplace} would produce a
holomorphic sum on a sector of opening greater than $\pi$ centered at the
negative axis.  By Section~\ref{sec:borel}, this function agrees with
$Q$ on a negative real interval.  The identity theorem makes it an
analytic continuation of $Q$ on the connected overlap with the left
half-disc.  The larger sector contains sufficiently small nonzero
points of both imaginary rays, with neighborhoods crossing the imaginary
axis.  This contradicts Section~\ref{sec:natural}.

Thus no such exponential bound exists.  Since $\B\Q$ is bounded on
compact subsets of the plane, the unbounded supremum in
\eqref{eq:no-angular-bound} must be witnessed at unbounded radii.
This proves both the growth assertion and failure of angular summability
for $\Q$.

Angularly $1$-summable tangent-to-identity series form a group under
composition, and their summation commutes with inversion; this is the
standard nonlinear inversion theorem, in small-variable notation
\cite[Theorem 17.3]{sauzin}.  Equivalently, it follows by changing from
$u=0$ to $z=\infty$, applying the tangent-to-identity inversion theorem
there, and using the algebra property for reciprocals.
If $\U$ were angularly summable in direction $\pi$, its formal inverse
$\Q$ would be so after shrinking the arc.  This contradiction proves
the assertion for $\U$.
\end{proof}

\begin{corollary}[Divergence despite an entire Borel transform]
\label{cor:divergence}
The series $\Q$ has radius of convergence zero, although
$\limsup(|[u^n]\Q|/n!)^{1/n}=0$.
\end{corollary}
\begin{proof}
Entirety of $\B\Q$ gives the displayed zero limsup.  A convergent power
series has a Borel transform of finite exponential type in every
direction, so would contradict Theorem~\ref{thm:main}(iii).
\end{proof}

\begin{remark}[Finite-plane resurgence versus infinity]
No finite Borel singularity can explain the obstruction in this example:
there are none for $\B\Q$.  In a definition of resurgence based only on
endless analytic continuation, the Borel transform is automatically
resurgent.  Assertions that it is ``not resurgent'' would therefore be
incorrect without an additional growth convention.  The precise
conclusion is failure of angular $1$-summability at $\pi$.
\end{remark}

\section{Rational amplitudes and robustness}
\label{sec:rational}

\subsection{A general quadratic-growth fine-summation lemma}
The Borel estimate extends beyond rational amplitudes.  Rationality will
be used only for the natural-boundary obstruction.

\begin{lemma}[Exponentially bounded amplitudes and quadratic slopes]
\label{lem:weighted-borel}
Suppose $c_1\ne0$, $|c_j|\le M D^j$, and the real slopes satisfy
\[
 -L\le\lambda_j\le L(1+j^2)
\]
for some positive $M,D,L$.  Then the small left-side inverse of
\eqref{eq:weighted-kernel} has an entire integrated Borel transform,
is finely summable in direction $\pi$, and is the corresponding fine
sum.  It is holomorphic in a small left half-disc and smooth on its
diameter.
\end{lemma}

\begin{proof}
Choose a $q$-disc on which the geometric majorants and all their
polynomially weighted derivatives converge.  For sufficiently small
$u$ in the closed left half-plane, solve the equation by isolating
$c_1\ee^{\lambda_1u}q$ and applying a uniform contraction to the
remaining terms.  Because $\lambda_j\ge-L$,
$|\ee^{\lambda_j u}|\le\ee^{L|u|}$ there.  The derivative at $q=0$
tends to $c_1\ne0$.  This proves existence, uniqueness, holomorphy and,
by the derivative argument of Section~\ref{sec:inverse}, smoothness.

The Lagrange blocks are indexed as before.  Their scalar weights and
frequencies are
\begin{align}
 w_{k,\mathbf m}
   &=\frac{(-1)^\ell}{k}\binom{k+\ell-1}{\ell}
       c_1^{-k-\ell}\prod_{i=1}^{\ell}c_{m_i+1},
       \label{eq:general-weight}\\
 b_{k,\mathbf m}
   &=\sum_{i=1}^{\ell}\lambda_{m_i+1}-(k+\ell)\lambda_1.
       \label{eq:general-frequency}
\end{align}
The initial block is $c_1^{-1}u\ee^{-\lambda_1u}$.
The exponential coefficient bound, the number of compositions, and the
binomial bound imply that the uncancelled total absolute weight of
block $k$ is at most $C_0^k$ for a constant $C_0$.
The slope assumptions give, after enlarging $C_1$,
\begin{equation}
 -C_1k\le b_{k,\mathbf m}\le C_1k^2.
 \label{eq:general-frequency-bounds}
\end{equation}
For the upper bound use
$\sum_i(m_i+1)^2\le(\sum_i(m_i+1))^2\le4k^2$;
for the lower bound use $\ell\le k$ and the uniform lower bound
on the slopes.

On any compact Borel disc the bounds from
\eqref{eq:bessel-integral} are $|\zeta|^k\exp(O(k))/k!$,
so the weighted Borel series converges normally and is entire.
For nonnegative frequencies in $T_\delta$, its terms are bounded by
\[
 C_0^k\frac{|\zeta|^k}{k!}
          \exp(2\sqrt{C_1\delta}\,k).
\]
For negative frequencies, use
\[
 2\sqrt{C_1k|\zeta|}\le k+C_1|\zeta|
\]
to obtain a bound
$C_0^k|\zeta|^k\ee^{k+C_1|\zeta|}/k!$.
Summing over $k$ yields an exponential bound on each fixed tube.

On the negative ray, the same estimates dominate the complete Borel
sum by $C\ee^{Hs}$ for some $C,H$.  Equation~\eqref{eq:block-laplace}
and dominated convergence identify its Laplace transform with the
convergent exponential-block inverse for small negative inputs.
Its equality with the literal branch follows from the parameterized
Lagrange formula and uniqueness.  This proves the lemma without
assuming positivity of the amplitudes.
\end{proof}

\subsection{Periodic twisting preserves rationality and infinite support}
Let $\rationalamp$ be as in \eqref{eq:amplitudes}, and let $f_j$ be a
nonzero periodic sequence whose values are all nonzero.  Then
\begin{equation}
 \rationalamp_f(z)=\sum_{j\ge1}c_jf_jz^j
 \label{eq:twisted-amplitude}
\end{equation}
is rational: expand $f_j$ in its finite Fourier series and write
$\rationalamp_f$ as a finite linear combination of
$\rationalamp(\omega z)$, where $\omega$ ranges over roots of unity.
If $\rationalamp_f$ were a polynomial, its coefficients would vanish
for all large $j$.  Since $f_j\ne0$, the same would hold for $c_j$,
so $\rationalamp$ would be a polynomial.  Thus a nonpolynomial rational
amplitude remains nonpolynomial after quadratic rational-phase twisting.

All finite poles of $\rationalamp_f$ lie at rotated poles of
$\rationalamp$; cancellations can remove poles but cannot introduce
smaller-modulus ones.  If $R_A$ is the radius of convergence of
$\rationalamp$, then for $|q_0|<R_A$, every pole of
$\phi(a)=\rationalamp_f(\ee^{-a})$ lies strictly to the left of
$a_0=-\log q_0$ in real part.  Indeed a pole $z_*$ gives
$\Ree\omega=-\log|z_*|$, whereas
$\Ree a_0=-\log|q_0|>-\log R_A$.
The poles form a finite union of vertical lattices and have finite orders.
Theorem~\ref{thm:rigidity} applies, including possible higher-order poles.

\subsection{Eventually quadratic slopes}
Assume now \eqref{eq:eventual-quadratic} and suppose
$\rationalamp$ is nonpolynomial.  The assumptions of
Lemma~\ref{lem:weighted-borel} hold, since the coefficients of any
rational germ are exponentially bounded and the slopes are bounded
below and grow at most quadratically.

Take a nonzero imaginary point $u_0$ sufficiently close to zero with
\begin{equation}
 \alpha u_0=\frac{2\pi i r}{m}.
 \label{eq:scaled-rational-time}
\end{equation}
Suppose that $Q$ extends analytically across $u_0$.  Set
$u=u_0+t/\alpha$.  Apart from finitely many terms, the left side of
\eqref{eq:weighted-kernel} is
\begin{equation}
 \ee^{\gamma u}\sum_{j\ge1}c_j
      \ee^{\alpha u_0j^2}
      \left(Q(u)\ee^{\beta u}\right)^j\ee^{j^2t}.
 \label{eq:weighted-heat-reduction}
\end{equation}
The discrepancy from the exceptional slopes is a finite analytic sum
under the assumed analyticity of $Q$.
The factor $Q(u_0)$ is nonzero, because the right side of
\eqref{eq:weighted-kernel} is $u_0\ne0$.
Choose the half-disc small enough that
$|Q(u_0)\ee^{\beta u_0}|<R_A$; this is possible because
$Q(u)\to0$ and $\beta$ is real.
An analytic logarithm of $Q(u)\ee^{\beta u}$ supplies an analytic curve
$a(t)$ in the heat reduction.

The twisted amplitude is rational and nonpolynomial by the preceding
subsection.  Its nearest heat poles lie strictly to the left of $a(0)$.
Theorem~\ref{thm:rigidity} shows that the heat expression in
\eqref{eq:weighted-heat-reduction} cannot be analytic in $t$.
But the defining equation, after subtracting the finite exceptional
terms and dividing by $\ee^{\gamma u}$, says that it is analytic.
This contradiction proves noncontinuability at every point
\eqref{eq:scaled-rational-time}.  Density gives the entire small
imaginary interval as a natural boundary.

\begin{proof}[Completion of Theorem~\ref{thm:rational-main}]
The preceding arguments prove the existence, Borel and natural-boundary
assertions in the nonpolynomial case.  Fine-sum identification and
Lemma~\ref{lem:sector-laplace} exclude angular summability just as in
the quadratic unit-amplitude model.  Positive Taylor radius would imply
angular summability and is therefore excluded as well.

If $\rationalamp$ is a polynomial, the kernel in
\eqref{eq:weighted-kernel} is a finite sum of entire functions of $u$
and polynomials in $q$.  It is holomorphic near $(0,0)$, with
$q$-derivative $c_1\ne0$.  The holomorphic implicit function theorem
gives a convergent inverse germ.  Its Borel transform has finite
exponential type in every direction and its usual sum agrees with the
fine sum.  All four stated conditions then hold.
\end{proof}

\section{Finite-action approximation does not detect the boundary}
\label{sec:truncation}

The natural boundary coexists with exceptionally strong approximation by
analytic finite-action models.  This is relevant to symbolic algorithms
and to interpreting finite numerical evidence.

For the unit quadratic model put
\[
 K_J(q,u)=\sum_{j=1}^Jq^j\ee^{j^2u},
 \qquad K_J(Q_J(u),u)=u.
\]
The same contraction used for $Q$ defines $Q_J$ on the same closed
half-disc.  Since the finite kernel is holomorphic across its diameter
and its $q$-derivative remains nonzero there, each $Q_J$ extends
holomorphically across every point of that diameter.  The size of the
extension neighborhoods is not uniform in $J$.

\begin{theorem}[Exponential approximation and exact finite jets]
\label{thm:truncation}
For every $J\ge1$,
\begin{equation}
 \sup_{\substack{|u|\le1/32\cr\Ree u\le0}}
   |Q(u)-Q_J(u)|
       \le\frac{7680}{5983}\,16^{-(J+1)}.
 \label{eq:truncation-bound}
\end{equation}
For each fixed integer $m\ge0$ there is a constant $C_m$, independent
of $J$, such that
\begin{equation}
 \max_{0\le h\le m}\ 
 \sup_{\substack{|u|\le1/32\cr\Ree u\le0}}
       |Q^{(h)}(u)-Q_J^{(h)}(u)|
       \le C_m(J+1)^{2m}16^{-(J+1)}.
 \label{eq:derivative-truncation}
\end{equation}
Furthermore,
\begin{equation}
 [u^n]\widehat Q_J=[u^n]\Q,\qquad n\le J.
 \label{eq:jet-equality}
\end{equation}
\end{theorem}

\begin{proof}
The difference between the full and truncated contraction maps is at most
\[
 \frac{32}{31}\sum_{j>J}(1/16)^j
       =\frac{32}{31}\frac{16^{-(J+1)}}{1-1/16}.
\]
Divide by $1-32/225$, the fixed-point stability denominator from
\eqref{eq:contraction-constant}.  The resulting constant is
$7680/5983$, proving \eqref{eq:truncation-bound}.

For derivatives, every mixed kernel derivative of total order at most
$m$ introduces at most a constant times $j^{2m}$ into the tail, up to a
fixed factor from $q$-differentiation.  Consequently its tail is bounded by
$C_m'(J+1)^{2m}16^{-(J+1)}$.  To compare evaluation at $Q$ and $Q_J$,
use the mean-value estimate in the $q$ variable; all needed kernel
$q$-derivatives are uniformly bounded on a disc containing both values,
because both actually lie in $|q|\le17/465<1/16$.
Repeated implicit differentiation expresses the $h$th derivative of
$Q$ as a rational polynomial in mixed kernel derivatives of order at
most $h$, the lower derivatives of $Q$, and $1/K_q$.
The denominators stay uniformly away from zero for all $J$.
Induction on $h$, using \eqref{eq:truncation-bound} at the initial step,
proves \eqref{eq:derivative-truncation}.  Additional $q$ derivatives used
only to compare evaluation points are uniformly bounded and multiply
the already controlled difference $Q-Q_J$; they do not increase the
power of $J$ in the tail bound.

Finally, the omitted kernel terms have $q$-order at least $J+1$.
Since the inverse starts with $u$, they contribute only from $u$-order
$J+1$ onward in the formal implicit equation.  Recursive uniqueness of
the coefficients proves \eqref{eq:jet-equality}.
\end{proof}

This theorem rules out two tempting but invalid inference patterns.
Convergence of analytic approximants in every fixed differentiability
norm on a one-sided domain does not supply analytic continuation of the
limit.  Likewise, agreement of arbitrarily long finite Taylor prefixes
with convergent finite-action models does not supply a common convergence
radius or angular Borel bound.

\section{Examples, checks, and limitations}
\label{sec:examples}

\subsection{A nonlinear analytic curve and one dominant pole}
Take
\[
 a(t)=\log2+\frac3{10}t+\frac17t^2,
 \qquad
 G(t)=\sum_{j\ge1}2^{-j}
          \exp\left(tj^2-\frac3{10}jt-\frac17jt^2\right).
\]
The initial function is $\phi(a)=1/(\ee^a-1)$.
The nearest pole to $a_0=\log2$ is the simple pole at zero with residue
one.  Theorem~\ref{thm:rigidity} gives
\begin{equation}
 \frac{g_n(\log2)^{2n+1}}{D_{1,n}}
       \longrightarrow2^{-3/20}.
 \label{eq:heat-check-limit}
\end{equation}
All $g_n$ are rational: the moments
$\sum_{j\ge1}j^h2^{-j}$ are integers for $h\ge1$, and the curve
coefficients are rational.  The accompanying program computes $g_n$
exactly before evaluating the normalization at 120-digit precision.

\begin{center}
% Generated by verify.py. Floating values are diagnostics, not interval certificates.
\begin{tabular}{rcc}
\toprule
$n$ & $g_n(\log 2)^{2n+1}/D_{1,n}$ & ratio to $2^{-3/20}$ \\
\midrule
10 & 0.900426863367 & 0.999086159422 \\
20 & 0.900850024301 & 0.999555685875 \\
40 & 0.901052952323 & 0.999780848615 \\
80 & 0.901152369379 & 0.999891158745 \\
\bottomrule
\end{tabular}

\end{center}

These values test the nontrivial multiplier, including a genuinely
nonlinear curve.  The quadratic coefficient of the curve changes finite
orders but not the limit.  The computation is not used in the proof.

\subsection{Two equally near poles: do not divide by an oscillating leading term}
Replace the amplitude $2^{-j}$ by $(-1/2)^j$, keeping the same curve
increment.  Now $\phi(a)=-1/(\ee^a+1)$ and the nearest poles to
$\log2$ are $\pm i\pi$, both with residue one.  Let
\[
 R=\sqrt{(\log2)^2+\pi^2},\qquad s_\pm=\log2\pm i\pi.
\]
The theorem gives the additive asymptotic
\begin{equation}
 \frac{g_nR^{2n}}{D_{1,n}}
   -\sum_{\sigma\in\{+,-\}}
      \frac{\ee^{-3s_\sigma/20}}{s_\sigma}
                (R/s_\sigma)^{2n}\longrightarrow0.
 \label{eq:two-pole-check}
\end{equation}
The leading sum oscillates and can be small.  A relative error would be
an unjustified strengthening.

\begin{center}\small
\begin{tabular}{rrrr}
\toprule
$n$ & normalized coefficient & nearest-pole model & difference \\
\midrule
10 & -0.4906207135 & -0.454908339 & -0.035712374 \\
20 & 0.4717175341 & 0.4692056033 & 0.0025119308 \\
40 & -0.5598737788 & -0.553093519 & -0.0067802598 \\
80 & 0.03428136919 & 0.03726718722 & -0.002985818 \\
\bottomrule
\end{tabular}

\end{center}

The nonmonotone displayed differences are compatible with an additive
$o(1)$ statement.  The proof uses the positive Ces\`aro mean square
\eqref{eq:cesaro}, not a pointwise lower bound at every index.

\subsection{An actual rational boundary point}
The point $u_0=2\pi i/257$ lies in the prescribed half-disc diameter.
With $J=96$, fixed-point evaluation gives
\[
 Q_J(u_0)\approx
 0.00119519260283291956428615710818
 +0.0244408539831444770352844651668\,i.
\]
The analytic action-truncation bound
\eqref{eq:truncation-bound} is less than
$2.04\times10^{-117}$.  This bounds the difference of the exact full
and exact finite-action functions; it does \emph{not} by itself include
floating-point rounding error.  The recorded 120-digit arithmetic is
therefore labeled a diagnostic, not an outward-rounded enclosure.
The theorem proves noncontinuability at this point from finite Fourier
pole data, without using the numerical value of $Q(u_0)$.

\subsection{Rational amplitudes covered by the dichotomy}
The amplitudes $c_j=1$, $c_j=j$, and $c_j=j^d$ for a fixed nonnegative
integer $d$ all have rational generating functions with infinite support.
So do nonzero periodic amplitudes after any finite initial correction
that preserves $c_1\ne0$.  Each is covered by
Theorem~\ref{thm:rational-main}.  Higher-order poles in the polynomial
amplitude examples are the reason the heat-rigidity theorem was proved
for arbitrary pole order rather than only simple poles.

Geometric amplitudes $c_j=\rho^j$, $\rho\ne0$, merely rescale the
$q$ variable and also lie in the class.  Finite changes of coefficients
or slopes cannot eliminate the obstruction as long as the amplitude
generating function remains nonpolynomial and the linear coefficient
remains nonzero.

\subsection{Why the rationality restriction cannot be dropped}
Consider $c_j=\ee^{-j^3}$ and $\lambda_j=j^2$.
On every compact set in $(q,u)$,
\[
 |c_jq^j\ee^{j^2u}|
       \le\exp(-j^3+Cj^2+C'j),
\]
so the kernel converges normally on all of $\C^2$.
Its derivative with respect to $q$ at $(0,0)$ is $\ee^{-1}\ne0$,
and its inverse is an ordinary holomorphic germ at zero.
This infinite-support example is outside the rational-amplitude class.
It shows that the general slogan ``infinitely many quadratic actions
force a natural boundary'' is false.  The amplitude hypothesis must
remain in the theorem.

\subsection{What the exact verification establishes}
The standard-library part of \texttt{verify.py} computes the formal
fixed point and its inverse independently of the signed-block formula.
It checks both compositions, the literal implicit coefficient residual,
the block expansion, the frequency bounds, the first Schr\"oder masses,
finite-action prefix agreement, and the rational contraction constants.
The recorded run checks all these identities through degree sixteen.
No finite verification is represented as a proof of the infinite
natural-boundary or summability statements.

\section{Further research questions}
\label{sec:future}

\subsection{Individual Borel rays and their exponential types}
The negative open-arc question is now settled, but the set
\[
 E_Q=\left\{\theta:
   \limsup_{r\to\infty}\frac1r
       \log^+|\B\Q(r\ee^{i\theta})|<\infty\right\}
\]
remains to be determined.  Theorem~\ref{thm:main} excludes a uniform
angular neighborhood at $\pi$; it does not say that every neighboring
individual ray fails.  A useful next result would distinguish isolated
finite-type rays from families with locally bounded type.

\subsection{The largest shrinking neighborhoods of the negative ray}
Fixed-width neighborhoods admit exponential bounds, whereas fixed-angle
wedges do not.  Determine a sharp width function $w(r)$ for domains
$|\Imm\zeta|\le w(-\Ree\zeta)$ far along the negative axis.
The Bessel estimate and
$|\Ree\sqrt\zeta|\asymp |\Imm\zeta|/\sqrt{|\Ree\zeta|}$ suggest
that parabolic and subparabolic scales are natural, but the uncancelled
bound is not an optimality proof.  This is a quantitative question about
cancellation at Borel infinity.

\subsection{The actual boundary jets of the nonlinear inverse}
The rigidity theorem tests a hypothetical analytic amplitude; it does
not compute the large derivatives of the actual nonanalytic $Q$ at a
rational boundary point.  Determine their exact Gevrey type, dominant
phases and possible transseries.  A successful calculation must solve a
nonlinear derivative recurrence rather than insert the nonexistent
analytic Taylor germ of $Q$ into
\eqref{eq:curved-pole-asymptotic}.

\subsection{Irrational boundary points and arithmetic dependence}
Density proves noncontinuability at irrational points without analyzing
their local coefficient growth.  Do Diophantine approximation properties
of $\Imm u/(2\pi)$ affect the optimal derivative bounds, quasi-analytic
classes, or local summation procedures there?  Rational approximants give
increasing periods, so uniform control in the period is the first
technical target.

\subsection{Where the finite-action singularities accumulate}
Every $Q_J$ extends across the diameter, yet the limit does not.
Locate the nearest branch points or other inverse singularities of
$Q_J$, determine their distance to the imaginary axis as $J\to\infty$,
and test whether their distribution converges to a measure on the
limiting boundary.  This would explain how a natural boundary emerges
from finite analytic models rather than merely proving that it does.

\subsection{Beyond rational amplitudes}
Classify amplitude generating functions for which analytic-curve
rigidity still forces a natural boundary.  Algebraic functions with
nonintegral local exponents are a natural next class: a corresponding
transfer lemma should replace the integer pole order by a complex local
exponent.  Entire amplitudes require a separate growth analysis, as the
$\ee^{-j^3}$ example demonstrates.  The desired criterion must combine
singularities of the amplitude generating function with the quadratic
slope scale.

\subsection{Higher-degree and nonpolynomial slopes}
For $\lambda_j\sim j^d$, $d>2$, the local formal operator is a
higher-order heat operator and the natural Gevrey scale changes.
Analytic-curve multipliers and cancellation among equal $d$th powers of
pole distances require new analysis.  The simple right-half-plane
argument excluding opposite squared distances is special to $d=2$.
For noninteger powers, rational boundary phases no longer become
periodic by the argument used here.

\subsection{Acceleration compatible with the literal branch}
The failure of an angular $1$-sum need not prohibit acceleration or a
summation scheme adapted to shrinking domains.  Construct such a
procedure from the formal coefficients, prove that it selects the
literal left-side inverse, and state exactly what continuation or
uniqueness property it supplies.  An acceleration theorem should not
silently claim analytic continuation through the proved natural
boundary.

\subsection{Formalization and finite certificates}
A useful formalization starts with the finite coefficient identities,
then proves the uniform domination in
\eqref{eq:summable-majorant}, the pole decomposition, and the Ces\`aro
noncancellation lemma.  These form a reusable analytic-curve rigidity
module.  A separate complex-analysis layer should establish smooth
implicit branches, boundary noncontinuation, and Laplace-sector gluing.
The existing repository infrastructure does not make those new analytic
obligations automatic.

\section{Conclusion}
The negative-ray summation problem and the angular summation problem
have different answers for quadratic exponential feedback.  The inverse
has an entire Borel transform, and its Laplace integral is controlled in
fixed-width Borel neighborhoods.  Nevertheless a smooth natural boundary
prevents every angular enlargement around the negative direction.
The obstruction survives finite coefficient and slope perturbations and,
more generally, gives a finite-support versus natural-boundary dichotomy
for rational amplitudes with eventually quadratic slopes.

The mechanism is not merely that a partial theta kernel is singular.
Its singularity cannot be removed by making its amplitude an analytic
unknown: each nearest heat pole retains a nonzero analytic-curve
multiplier, and the finite pole contributions cannot cancel permanently.
This closes the specific implicit-cancellation gap required to answer
the repository's open-arc question.

\appendix
\section{Proof-dependency and claim audit}
\label{app:audit}

The natural-boundary theorem for the unit model depends only on the
uniform implicit construction, the factorial composition lemma, analytic
remainder majorization, the finite-pole Ces\`aro argument, and elementary
periodic Fourier inversion.  It does not use the earlier regularity
article's large-order asymptotics, a general transseries realization
theorem, or a numerical computation.

The entire-Borel and fine-summation statements are proved directly from
Lagrange inversion and the normalized Bessel integral.  The block
representation is credited to the earlier repository work in its general
form.  The angular obstruction for $\Q$ then uses only elementary
Laplace contour deformation and the identity theorem.  The extension of
that obstruction to the formal forward series $\U$ uses the standard
angular summability inversion theorem \cite{sauzin}.  The broader
rational-amplitude classification uses the same elementary proof chain
with a weighted block estimate and rational periodic twisting.

The nowhere-real-analytic assertion concerns the boundary \emph{curve},
not just one parametrization, and uses local analytic straightening and
Schwarz reflection.  The finite-action theorem proves convergence on a
closed one-sided domain; it does not supply a common two-sided
holomorphic neighborhood.  The numerical tables are diagnostics, not
interval-certified real inequalities or substitutes for the proofs.

No Lean build was attempted and no Lean proof file is included.  No
repository branch or file was modified.  The inspected research packages
are themselves marked unmerged and not formalized in the repository
README \cite{repo-index}.  The present article should receive an
independent mathematical review before being treated as an audited result.

\section{Reproducibility and source provenance}
\label{app:reproduce}

The package contains the editable \texttt{article.tex}, the compiled
\texttt{article.pdf}, \texttt{verify.py}, generated exact and numerical
JSON results, the two small TeX table inputs, the captured verification
output, a README, and a source-provenance note.
The TeX does not depend on a local copy of ProveIt or on custom fonts.

The commands used are
\begin{quote}\small\ttfamily
python verify.py\\
pdflatex -interaction=nonstopmode -halt-on-error article.tex\\
pdflatex -interaction=nonstopmode -halt-on-error article.tex\\
pdflatex -interaction=nonstopmode -halt-on-error article.tex
\end{quote}
The exact checks use Python's \texttt{fractions.Fraction}.  The numerical
diagnostics use \texttt{mpmath} 1.3.0 at 120 decimal digits.  The program
makes no network requests.  Its order can be increased up to twenty with
\texttt{--max-order}; the independent composition enumeration has
exponential cost and is intentionally bounded.

The following repository files were directly read at the pinned commit:
\begin{itemize}
\item \path{Analysis/Transseries/docs/series-and-transseries/README.md},
including its descriptions of the recently filed feedback articles;
\item \path{Negative_Ray_Summation_Exponential_Feedback/article.tex}
under that directory, including its definitions, Lagrange blocks,
Bessel representation and explicit angular-summability question.
\end{itemize}
The source review was focused, not an exhaustive audit of every statement
in the repository's large transseries volumes.  The reference to the
regularity article supplies context as recorded in the directly read
index and negative-ray paper; the present proof does not depend on its
sharp asymptotic theorem.

\begin{thebibliography}{9}
\raggedright

\bibitem{repo-index}
V. Reshetnikov, \emph{ProveIt: series and transseries}, repository index,
commit \texttt{0c973d8f5e7ef4550f4df1bf486d890037fd9f14}, accessed
29 September 2026.
\url{https://github.com/VladimirReshetnikov/ProveIt/blob/0c973d8f5e7ef4550f4df1bf486d890037fd9f14/Analysis/Transseries/docs/series-and-transseries/README.md}.

\bibitem{repo-negative}
\emph{Negative-Ray Summation of Countable Exponential-Feedback Transseries},
research draft prepared for V. Reshetnikov, 29 September 2026,
ProveIt, same pinned commit, directory
\path{Analysis/Transseries/docs/series-and-transseries/}
\path{Negative_Ray_Summation_Exponential_Feedback/},
\texttt{article.tex}.  In particular, the research-agenda subsection
``Angular summability of the quadratic model''; the direct source range
read for that question was lines 1330--1410.

\bibitem{repo-regularity}
\emph{A Sharp Regularity Classification for Countable Exponential-Feedback
Transseries}, research draft in ProveIt, directory
\path{Exponential_Feedback_Regularity_Classification/}
under the same series-and-transseries group, as described in
\cite{repo-index,repo-negative}.  Contextual reference; its sharp
coefficient asymptotics are not a dependency here.

\bibitem{sauzin}
D. Sauzin, \emph{Introduction to 1-summability and resurgence},
arXiv:1405.0356, 2014.  Sections 7 and 9 distinguish fine summability
from summability in an arc; Section 17 proves the nonlinear group and
inverse properties.
\url{https://arxiv.org/abs/1405.0356}.

\bibitem{lms}
D. A. Lutz, M. Miyake, and R. Sch\"afke,
On the Borel summability of divergent solutions of the heat equation,
\emph{Nagoya Mathematical Journal} \textbf{154} (1999), 1--29.
DOI: \href{https://doi.org/10.1017/S0027763000025289}{10.1017/S0027763000025289}.

\bibitem{mp}
S. Michalik and B. Podhajecka,
\emph{The Stokes phenomenon for certain partial differential equations
with meromorphic initial data}, arXiv:1603.04209, 2016.
\url{https://arxiv.org/abs/1603.04209}.

\bibitem{hlss}
L. Han, Y. Li, D. Sauzin, and S. Sun,
\emph{Resurgence and Partial Theta Series}, arXiv:2112.15223v3, 2022.
\url{https://arxiv.org/abs/2112.15223}.

\bibitem{gessel}
I. M. Gessel, Lagrange inversion,
\emph{Journal of Combinatorial Theory, Series A} \textbf{144} (2016),
212--249; arXiv:1609.05988.
\url{https://arxiv.org/abs/1609.05988}.

\end{thebibliography}
\end{document}
